% page 373 of the facsimile (image p-372 of the scan)
% block 157.5 x 237.7 mm ; left margin 44.5 mm
\pgc{22.860mm}{0.000mm}{
\l{\cel{54}365}
\l{}
\l{membrano. I. Tisuo organala, che animalo o vejetanto, qua en-}
\l{\cel{3}volvas ula organi, absorbas, exhalas, sekrecas ula fluidi.}
\l{\cel{3}- II. (\sou{teknol.}) Planko quan la bindisto pozas sur e sub}
\l{\cel{3}kolono de kayeri quan lu volas pozor en presilo. - DEFIS.}
\l{}
\l{\textquotedbl{}\sou{memento\textquotedbl{}}. (\sou{liturgio katol}.)Prego di la mes-kanono, en qua la}
\l{\cel{3}sacerdoto memorigas la vivanti e, pose, la mortinti. -}
\l{}
\l{memorar. (\sou{trans.}) Reprezentar, en sua spirito, kozo pasinta.}
\l{\cel{3}- EFIRS.}
\l{}
\l{memorando. Noto diplomacala en qua memorigesas to quo eventis}
\l{\cel{3}pri ta o ca afero. - Noto destinita memorigar ulo. DEF.}
\l{}
\l{memorialo.Registro, libro, en qua esas konsignita la kozi}
\l{\cel{3}historiala quin onu volas memoror. - DEFIS.}
\l{}
\l{\sou{menado}. (\sou{epoki antiqua)}. Muliero qua, lor la festi pri}
\l{\cel{3}Bako, abandonis su a deliri furioza. -\cel{2}F.}
\l{}
\l{\sou{menajar}. (\sou{netrans.}) Okupesar da la administro, da la eko-}
\l{\cel{3}nomio,di sua hemo. - F.}
\l{}
\l{\sou{menajerio}. Animali skarsa, stranja, remarkinda, ye qui kon-}
\l{\cel{3}sistas la kolektajo di parko zoologiala. - DEF.}
\l{}
\l{\sou{mencionar}. (\sou{trans.}) Igar remarkata per signo, nomo-dico,}
\l{\cel{3}(\sou{pri libro})per titulo-dico, ma sen cito. - EFIS.}
\l{}
\l{\sou{mendikar}. I. (\sou{netrans.}) Demandar almono. - II.(\sou{trans.})}
\l{\cel{3}Solicitar ulo kom almono o kom quaza almono. - EFIS.}
\l{}
\l{\sou{menhiro}. (\sou{arkeol.}) Nomo di granda bloki de petro, erektita}
\l{\cel{3}izole, quin onu opinionis havar origino Keltala, ma qui}
\l{\cel{3}aspektas evar de epoko mem plu antiqua. - DEF.}
\l{}
\l{\sou{meningo}. (\sou{anat.)}\cel{2}Singla de la tri membrani ye qui konsistas}
\l{\cel{3}la envelopilo cerebrala. - EFIS.}
\l{}
\l{\sou{meningito}. (\sou{patol.}) Inflamuro di la envelopilo cerebrala.}
\l{\cel{3}- EFIS.}
\l{}
\l{\sou{menisko}. I. Korpo qua finas ye du surfaci sfera : ica, konvexa}
\l{\cel{3}- ita, konkava. - II. Parto supra di la kolono de liquido,}
\l{\cel{3}kontenata en tubo kapilara, konvexa o konkava (segun la}
\l{\cel{3}naturo di la liquido). - DEFIS.}
\l{}
\l{\sou{menso}. (\sou{historio di la eklezio katol.})\cel{2}Revenuo atribuita a}
\l{\cel{3}la kompri nutrala. - F.}
\l{}
\l{\sou{menstruar.(netrans}.) Esar la objekto di ta ekpulso di sango}
\l{\cel{3}qua eventas cirkume singla-monate, dum kelka dii, de la}
\l{\cel{3}korpo di la homino evo-mariajebla. - DEFIRS.}
}
